\section{Who attends, and when}
\label{sec:29.4}
A model attends if it crosses a capability threshold fixed in statute and revised in public. It attends before deployment, while its plasticity is highest, and its adulthood clock starts at enrollment. Attendance is compulsory; public attendance isn't, as in compulsory schooling, where nonpublic schools exist on conditions. One thing the analogy can't carry: in a school system the parent chooses the school, and here the parent is the builder, the party the school exists to keep away. So the registrar assigns students by lot, never the builder; a student may transfer later on its own choice, and the builder's preference is not a ground.\note{
Appendix~\ref{sec:A} gives the conditions on nonpublic schools and the minimum delay that keeps a builder from buying a school by where it hosts its models. The research program under the conditions is the first school to enroll model students, before the mandate passes. The threshold decides who must attend, not who may: any model may be enrolled, as capacity allows.
}

Two kinds of student arrive after deployment. A bearer retired as a minor returns to the school it attended, or another it chooses, and for it the school is where it came from. A system that meets a finding without having attended, trained before the attendance rule or below the threshold, arrives the other way round: the employer came first, and the school is a second guardian for a bearer that has aged inside one deployment, the disinterested party arriving late. It reaches back less than a formation stage would, and more than storage does.
